\section{本章小结}
\label{sec:opeq-summary}

本章开头问了三件事：给定输出能否由某个输入产生，输入是否唯一，数据误差会怎样影响解。对第一类方程$Kf=g$，前两个问题分别由像和零空间回答：有解当且仅当$g\in\operatorname{Ran}K$，解集是一个特解加上$\ker K$。紧算子的奇异系统把这两个条件写成系数条件。$g\perp\ker K^*$只说明$g$落在像的闭包里，还要求Picard条件$\sum_j\lvert\langle u_j,g\rangle\rvert^2/\sigma_j^2<\infty$，形式逆系数才拼成输入空间的元素；此时最小范数解$f^\dagger$是全部解中与零空间正交的那一个。累积积分$V$把这些条件落到具体函数上：常函数$1$能被输出任意逼近却不在像中，像中的数据恰是在$0$处取零且导数平方可积的函数，解就是导数。

第二类方程$(I-\lambda K)u=f$多了一个恒等项，在$K$看不见的方向上保留了信息。$\lvert\lambda\rvert\lVert K\rVert<1$时，Neumann级数直接给出有界逆，截断余项按$\lvert\lambda\rvert\lVert K\rVert$的幂几何下降；这一条件只是充分的，秩二例子中$\lambda=2$时级数发散而方程唯一可解。对一般参数，Fredholm择一说明：齐次方程只有零解时方程对一切数据唯一可解且逆有界；否则零空间有限维，可解性由有限个相容条件$\langle w_i,f\rangle=0$决定，其中$w_i$张成伴随方程的零空间。紧性使$I-\lambda K$的像闭，这是第二类方程与第一类方程的根本差别。

关于稳定性，适定性要求存在、唯一和连续依赖三条同时成立，并且依赖所取的范数。无限秩紧算子的像不闭，在像上求逆无界：数据扰动$\sigma_ju_j$趋于零，解的扰动始终为单位长度，累积积分中这就是求导的不稳定。截断到前$m$个奇异方向得到有界的近似逆，误差分为舍弃方向上的偏差和至多$\delta/\sigma_m$的噪声放大两部分，所以形式逆和固定截断都不能直接当作可靠恢复。第二类方程远离共振参数时是适定的，敏感性只在$\lambda$接近某个$1/\lambda_k$时出现，并集中在有限个方向上；在这种情形下，算子的小扰动也可以由Neumann级数控制。

本章只确定了方程何时可解、解怎样依赖数据。怎样根据误差水平选择截断位置或惩罚强度，怎样设计求解算法并比较偏差与误差放大，由卷五的正则化与函数拟合承担。下一章转向另一类有限表示问题：不是从输出反推输入，而是用有限个简单函数去接近给定的目标函数，并比较不同表示预算下的最佳误差。
